% ============================================================
% info.tex — 本科生毕业论文（设计）基本信息（SSOT）
% ============================================================
% 全篇所有需要个人填写的字段都集中在这一个文件。改一处，封面、
% 页眉、致谢所有相关位置自动同步。其余章节文件 (sections/*.tex)
% 只写论文内容，不需要再放姓名、日期等元数据。
%
% 编辑顺序建议：
%   ① 第 1 区：填封面字段
%   ② 第 2 区：注册正文反复出现的占位串（研究对象、数据集名、方法名…）
% ============================================================

% ---------- 1. 封面字段 ----------
\thesistitle{【中文论文题目：请替换为最终题目】}
\thesistitleen{【English Thesis Title: Replace with the Final Title】}
\thesisauthor{【学生姓名】}
\advisor{【导师姓名】\quad 【职称】}
\coadvisor{}                                                % 合作指导教师，无则留空
\major{【专业名称】}
\college{【学院名称】}
\thesisdate{【20XX 年 X 月】}

% ---------- 2. 正文可复用变量 ----------
% 在正文里用 \thevar{key} 引用；改这里全文统一。
% 未注册的 key 编译时给警告，便于排查。
\thesisvar{研究对象}{【研究对象，如：冬小麦】}
\thesisvar{数据集}{【数据集名称，如：Sentinel-2 时序】}
\thesisvar{方法名}{【本文方法的简称】}
\thesisvar{研究区}{【研究区或实验平台】}
